\documentclass[border=8pt]{standalone}
\usepackage{circuitikz}
\usetikzlibrary{positioning}
\ctikzset{resistors/scale=0.65, amplifiers/scale=0.8}
\begin{document}
\begin{circuitikz}[font=\footnotesize]

  % ======== stage 1: difference amplifier forms the error ========
  \node[op amp] (A1) at (3.4,0) {};
  \node[above=0.9cm of A1, font=\footnotesize\bfseries] {Stage 1: error};
  % inverting input from the feedback
  \draw (0,-0.3) node[left]{$V_{fb}$} to[R=$10\mathrm{k}$] (2.2,-0.3) -- (A1.-);
  % feedback resistor
  \draw (A1.-) -- ++(0,0.95) coordinate (f1) to[R=$10\mathrm{k}$] (f1-|A1.out) -- (A1.out);
  % non-inverting input from the setpoint, then to ground
  \draw (0,-1.5) node[left]{$V_{ref}$} to[R=$10\mathrm{k}$] (2.2,-1.5) -- (A1.+);
  \draw (A1.+) -- ++(0,-0.95) coordinate (g1) to[R=$10\mathrm{k}$] ++(0,-1.1) node[ground]{};
  \draw (A1.out) -- ++(0.5,0) coordinate (e) node[above right=-1pt]{$V_e$};

  % ======== stage 2: inverting attenuator sets Kp ========
  \node[op amp] (A2) at (7.8,-0.3) {};
  \node[above=0.9cm of A2, font=\footnotesize\bfseries] {Stage 2: gain $K_p$};
  \draw (e) -- ++(0,-0.3) coordinate (e2) to[R=$100\mathrm{k}$] (6.6,-0.6) -- (A2.-);
  \draw (A2.-) -- ++(0,0.95) coordinate (f2) to[R=$500\,\Omega$] (f2-|A2.out) -- (A2.out);
  \draw (A2.+) -- ++(-0.5,0) -- ++(0,-0.6) node[ground]{};
  \draw (A2.out) -- ++(0.5,0) coordinate (g) node[above right=-1pt]{$-K_pV_e$};

  % ======== bias divider off the +15 V rail ========
  \draw (10.2,2.6) node[above]{$+15\,\mathrm{V}$} to[R=$4\mathrm{k}7+2\mathrm{k}$] (10.2,1.3)
        coordinate (bias) to[R=$1\mathrm{k}2$] (10.2,0.0) node[ground]{};
  \node[left=2pt] at (bias) {$V_{bias}$};

  % ======== stage 3: difference amplifier adds the hover bias ========
  \node[op amp] (A3) at (13.8,-0.3) {};
  \node[above=1.4cm of A3, font=\footnotesize\bfseries] {Stage 3: level shift};
  % inverting input takes the gain stage output
  \draw (g) -- ++(0,-0.3) coordinate (g3) to[R=$10\mathrm{k}$] (12.6,-0.6) -- (A3.-);
  \draw (A3.-) -- ++(0,0.95) coordinate (f3) to[R=$10\mathrm{k}$] (f3-|A3.out) -- (A3.out);
  % non-inverting input takes the bias
  \draw (bias) -- (11.0,1.3) -- (11.0,-1.9) to[R=$10\mathrm{k}$] (12.6,-1.9) -- (A3.+);
  \draw (A3.+) -- ++(0,-0.95) coordinate (g3b) to[R=$10\mathrm{k}$] ++(0,-1.1) node[ground]{};
  \draw (A3.out) -- ++(0.7,0) node[right, align=left]{$V_{cmd}=V_{bias}+K_pV_e$\\to ADC input 1};

\end{circuitikz}
\end{document}
